\section{Problem 3: Food Supply}
Three villages depend on a truck for food deliveries while consuming food continuously.
Determine how long they can survive with the original truck capacity and whether increasing its capacity allows them to be supplied indefinitely.
	% We dont really need to add the process, and how/what order we implemented it
	% This subsection is just that so I'll just comment it out
	% 
	% \subsection{Approach}
	% 
	% When I started I did not know how I was supposed to solve this. I could
	% not see a way to dynamically run more steps in z3. The trick is to
	% realize you do not need to. You can dynamically build a solver for $1,
	% 2, ... n$ until you find what you're looking for. When I figured that
	% out it wasn't that hard anymore.
	% 
	% I then created a lot of util functions and types to make implementing
	% the assertions easier. When I didn't really knew how to implement the
	% code I went back to the report to instead do it as logic formulas.
	% This was easier to implement then code because of the lack of overhead
	% of those very types.
	% 
	% When the logic was finished in the report, it was then really easy to
	% implement the same assertions as code using those helper types.
    \subsection{Solutions}
        
        \subsubsection{constants}
            \begin{itemize}
                \item truck capacity
                    $\texttt{trucks\_max\_capacity} = \begin{cases}
                        130 & \text{for sub-question a } \\
                        150 & \text{for sub-question b } \\
                    \end{cases}$
                \item The max steps to emulate:
                    \texttt{max\_steps}
            \end{itemize}

        \subsubsection{Variables}
            \begin{itemize}
                \item
                    $\mathit{f}_{s,v} \in \mathbb{Z}$ representing the food left in
                    village $v$, at step $s$.
                \item
                    $\beta_{s} \in \mathbb{Z}$ representing the food in the
                    truck at step $s$.
                \item
                    $\mathit{t}_{s} \in \mathbb{Z}$ representing the total
                    time passed at step $s$.
                \item
                    $\ell_{s,l} \in \mathbb{B}$ representing whether the
                    truck is at location $l$ on step $s$.
            \end{itemize}

        \subsubsection{Sets}
            \begin{itemize}
                \item 
                    The set of villages: $V = \big\{ A, B, C \big\}$.
                \item 
                    The set of supplies: $W = \big\{ S \big\}$.
                \item
                    The set of locations: $L = V \cup W$
                \item
                    The set of routes: $R = \big\{
                        (A, B), (B, A),
                        (A, S), (S, A),
                        (A, C), (C, A),
                        (B, C), (C, B),
                        (C, S), (S, C)
                    \big\}$
                \item
                    The set of steps: $\mathbb{S} = \big\{
                        1, 2, ..., \texttt{ max\_steps}
                    \big\}$
                \item
                    The set of linked steps: $\mathbb{S}_t = \big\{
                        (1, 2), (2, 3), ..., 
                        (\texttt{ max\_steps} - 1,\texttt{ max\_steps})
                    \big\}$
            \end{itemize}

        \subsubsection{Functions}        
            \begin{itemize}
                \item 
                    A function the capacity of a village $v$:
                    \[
                        c : V \to \mathbb{N},\qquad c(v) = \begin{cases}
                            90  & \text{if } v = A \\
                            120 & \text{if } v = B \\
                            90  & \text{if } v = C
                        \end{cases}
                    \]
                \item 
                    A function to get the distance between two locations:
                    \[
                        d : R \to \mathbb{N},\qquad d(r) = \begin{cases}
                            15 & \text{if } r \in \big\{
                                (A, S), (S, A), (C, S), (S, C)
                            \big\} \\
                            17 & \text{if } r \in \big\{
                                (A, B), (B, A)
                            \big\}  \\
                            13 & \text{if } r = (B, C) \\
                            9 & \text{if } r = (C, B) \\
                            12 & \text{if } r \in \big\{
                                (A, C), (C, A)
                            \big\}  \\
                            1 & \text{if } r \in \big\{
                                (S, S), (A, A), (B, B), (C, C)
                            \big\} \\
                        \end{cases}
                    \]
                    note that the distance to the same place is $1$. This is
                    done on purpose. This allows you to "travel" to the same
                    place in $1$ time unit. Effectively allowing you to wait.
            \end{itemize}
    
        \subsubsection{Assertions}
            \begin{itemize}
                \item 
                    Asserting that a village at any step has a valid amount of food:
                    \[
                        \bigwedge_{(s, v) \in \mathbb{S} \times V} \bigg(
                            0 < \mathit{f}_{s,v} \le c(v)
                        \bigg)
                    \]
                \item 
                    Asserting that the truck at any step has a valid amount of food:
                    \[
                        \bigwedge_{s \in \mathbb{S}} \bigg(
                            0 < \beta_{s} \le \texttt{trucks\_max\_capacity}
                        \bigg)
                    \]
                \item 
                    Asserting that the truck must be at exactly one location:
                    \[
                        \bigwedge_{s \in \mathbb{S}} \bigg(
                            1 = \sum_{l \in L} \Big(
                                \operatorname{ite}\big(
                                    \ell_{s,l},\ 1,\ 0
                                \big)
                            \Big)
                        \bigg)
                    \]
                \item 
                    Asserting that you can only travel routes:
                    \[
                        \bigwedge_{\substack{
                            (s_1, s_2) \in \mathbb{S}_t \\
                            (l_1, l_2) \in L \times L
                        }} \bigg(
                            \Big( 
                                    \ell_{s_1,l_1}
                                \wedge
                                    (l_1, l_2) \not \in R
                            \Big) \implies \neg \ell_{s_2,l_2}
                        \bigg)
                    \]
                \item 
                    Asserting that time passes as you travel:
                    \[
                        \bigwedge_{\substack{
                            (s_1, s_2) \in \mathbb{S}_t \\
                            (l_1, l_2) = r \in R
                        }} \bigg(
                            \Big( 
                                \ell_{s_1,l_1} \wedge \ell_{s_2,l_2}
                            \Big) \implies \Big( 
                                \mathit{t}_{s_2} = \mathit{t}_{s_1} + d(r)
                            \Big)
                        \bigg)
                    \]
                \item 
                    Asserting that the truck is refilled at a supply:
                    \[
                        \bigwedge_{\substack{
                            s \in \mathbb{S} \\
                            l \in W
                        }} \bigg(
                            \ell_{s,l} \implies \beta_{s} = \texttt{trucks\_max\_capacity}
                        \bigg)
                    \]
                \item 
                    Asserting that the food goes down as time passes, but if
                    the truck is at that location, that location may refill.
                    \[
                        \bigwedge_{\substack{
                            (s_1,s_2) \in \mathbb{S}_t \\
                            v \in V
                        }}
                        \left(
                            \mathit{f}_{s_2,v}
                            =
                            \mathit{f}_{s_1,v}
                            -
                            \Delta_\mathit{t}
                            +
                            \operatorname{ite}\big(
                                \ell_{s_2,v}, 
                                \Delta_\beta, 
                                0
                            \big)
                        \right)
                        \qquad \text{ where } 
                            \Delta_\mathit{t} = \mathit{t}_{s_2} - \mathit{t}_{s_1}
                            \text{ and }
                            \Delta_\beta = \beta_{s_1} - \beta_{s_2}
                    \]
                \item 
                    Asserting the starting state, all villages are full, and so is the truck:
                    \[
                        \beta_{s} = \texttt{trucks\_max\_capacity}    
                        \wedge 
                        \left(
                            \bigwedge_{\substack{
                                v \in V
                            }} \mathit{f}_{s,v} = c(v)
                        \right)
                        \qquad \text{ where } 
                            s = 0
                    \]
                \item 
                    Asserting that food in the truck can not go up on village tiles:
                    \[
                        \bigwedge_{\substack{
                            (s_1, s_2) \in \mathbb{S}_t \\
                            l \in L
                        }}\left(
                            \ell_{s_2,l} \wedge l \not\in W \implies \beta_{s_2} \le \beta_{s_1}
                        \right)
                    \]
            \end{itemize}

        \noindent For sub-question A we just wish to prove that we cannot keep these
        satisfying all these constraints. So we keep increasing the amount of steps
        until there is no way to satisfy them.

        For sub-question B we also wish to prove that we can keep going
        forever. We are going to do this by proving there is some order of steps
        such that we get to the exact same state again. Proving that we can keep
        going forever. For that we've made the following assertion:
        \[
            \left(
                \bigvee_{\substack{
                    (s_1, s_2) \in \mathbb{S} \times \mathbb{S}
                }}
                % is not the same time
                s_1 \not = s_2 
                % Truck has same amount of food
                \wedge \beta_{s_1} = \beta_{s_2}
                % Truck is in same location
                \wedge \left(
                    \bigwedge_{\substack{
                        l \in L
                    }} \ell_{s_1,l} = \ell_{s_1,l}
                \right)
                % all villages have the same amount of food as before
                \wedge \left(
                    \bigwedge_{\substack{
                        v \in V
                    }} f_{s_1,v} = f_{s_2,v}
                \right)
            \right)
        \]

    \subsection{Results}

        \subsubsection{Visualisation Exercise 3a}
            This was UNSAT.

        \subsubsection{Visualisation Exercise 3b}

            \begin{table}[h]
                \centering
                \label{tab:exercise3b}
                \begin{tabular}{cc|ccccc}
                    \textbf{$s$} & \textbf{$t_s$} & \textbf{$F_{A,s}$} &
                    \textbf{$F_{B,s}$} & \textbf{$F_{C,s}$} & \textbf{$\beta_s$} &
                    \textbf{$\ell_{s,?}$} \\
                    \hline
                    0 & 0   & 60 & 60  & 60 & 150 & S \\
                    1 & 15  & 67 & 45  & 45 & 128 & A \\ \rowcolor{gray!15}
                    2 & 30  & 52 & 30  & 30 & 150 & S \\
                    3 & 45  & 37 & 15  & 16 & 149 & C \\
                    4 & 54  & 28 & 100 & 7  & 55  & B \\
                    5 & 67  & 15 & 87  & 39 & 10  & C \\
                    6 & 82  & 0  & 72  & 24 & 150 & S \\
                    7 & 97  & 79 & 57  & 9  & 56  & A \\
                    8 & 109 & 67 & 45  & 45 & 8   & C \\ \rowcolor{gray!15}
                    9 & 124 & 52 & 30  & 30 & 150 & S \\
                \end{tabular}
                \caption{The steps and the results for 3b proving perpetual looping.}
            \end{table}
